\documentclass[11pt]{article}
% =========================================================
%  學生講義 共用前言（以 XeLaTeX 編譯）
%  需要套件：xeCJK, tcolorbox（其餘多在 BasicTeX 內）
% =========================================================
\usepackage[a4paper,margin=2.3cm]{geometry}
\usepackage{amsmath,amssymb}
\usepackage{xcolor}
\usepackage{graphicx}

% ---- 中文字型（macOS 系統字型）----
\usepackage{xeCJK}
\setCJKmainfont{Songti TC}      % 內文：宋體
\setCJKsansfont{Heiti TC}       % 標題/強調：黑體
\xeCJKsetup{AutoFakeBold=2.2, AutoFakeSlant=0.2}

% ---- 版面 ----
\setlength{\parindent}{0pt}
\setlength{\parskip}{0.55em}
\linespread{1.15}
\renewcommand{\baselinestretch}{1.15}

% ---- 顏色 ----
\definecolor{cBlue}{HTML}{1f6feb}
\definecolor{cGray}{HTML}{6b7280}
\definecolor{cGrayBg}{HTML}{f3f4f6}
\definecolor{cPurp}{HTML}{7a3ff2}
\definecolor{cPurpBg}{HTML}{f5f1ff}

% ---- 標註框 ----
\usepackage[most]{tcolorbox}

% 就地補充新概念（灰底）：\begin{補充框}{標題}
\newtcolorbox{補充框}[1]{
  enhanced, breakable, arc=2pt, boxrule=0.6pt,
  colback=cGrayBg, colframe=cGray,
  left=9pt,right=9pt,top=7pt,bottom=7pt,
  before upper={\textbf{\sffamily #1}\par\vspace{3pt}},
}

% 延伸 / 開放問題（紫色左邊條）：\begin{延伸框}{標題}
\newtcolorbox{延伸框}[1]{
  enhanced, breakable, arc=2pt, boxrule=0pt,
  colback=cPurpBg, colframe=cPurpBg,
  borderline west={3pt}{0pt}{cPurp},
  left=10pt,right=9pt,top=7pt,bottom=7pt,
  before upper={\textbf{\sffamily 延伸閱讀｜#1}\par\vspace{3pt}},
}

% ---- 圖：置中、底下小字說明 ----
% \fig{檔名}{寬度}{說明}
\newcommand{\fig}[3]{%
  \begin{center}
  \includegraphics[width=#2\linewidth]{#1}\\[2pt]
  {\small\color{cGray} #3}
  \end{center}}

% ---- 章節標題用黑體 ----
\newcommand{\h}[1]{\sffamily #1}


\begin{document}

\begin{center}
{\sffamily\Huge\bfseries 必物-6 量子現象}\\[3pt]
{\sffamily\large 普通高中 物理（必修）・學生講義}
\end{center}
\vspace{0.6em}

二十世紀初，物理學看起來幾乎「完工」了：牛頓力學、電磁學、熱力學三大支柱解釋了從行星到火車的一切。然而就在這時，出現了三個古典物理\textbf{完全算不對}的實驗——\textbf{黑體輻射}、\textbf{光電效應}、\textbf{原子光譜}。為了解釋它們，物理學家被迫接受一個違反直覺的新觀念：能量不是連續的，而是\textbf{一份一份}的，這就是「量子化（quantization）」。

由此延伸，人們發現光既是波也是粒子、電子既是粒子也是波——\textbf{波粒二象性（wave–particle duality）}；最後再用「能階」解釋了原子為何只發出固定幾種顏色的光。這一章是「古典物理」與「近代物理」的分水嶺，也是後續所有現代科技（雷射、半導體、電子顯微鏡）的根。本章重在\textbf{定性理解}，定量計算留待高三。

\medskip
本章主線：

\begin{center}
黑體輻射 $\rightarrow$ 能量量子化 $\rightarrow$ 光電效應（光的粒子性）$\rightarrow$ 德布羅意波（物質的波動性）$\rightarrow$ 雙狹縫與波粒二象性 $\rightarrow$ 原子能階與光譜
\end{center}

\section*{\sffamily 6-1\quad 量子論的誕生：黑體輻射}

\subsection*{\sffamily A.\ 一個尷尬的失敗}

任何溫度高於絕對零度的物體都會輻射電磁波（鐵會燒紅、燈絲會發光）。物理學家想用古典電磁學加熱力學，算出一個「黑體（blackbody）」在各個波長放出多少能量。

\begin{補充框}{先搞懂：什麼是「黑體」？它跟「黑色」有關嗎？}
\textbf{黑體}＝一個能把照到它身上的所有電磁波（各種波長）\textbf{全部吸收}、完全不反射也不穿透的理想物體。因為它一點光都不反射回來，常溫下看起來就是全黑的——名字便由此而來。

最容易誤會的是：\textbf{黑體同時也是「最會放出」輻射的物體}。

\begin{itemize}\setlength{\itemsep}{1pt}
\item \textbf{常溫}：黑體放出的是紅外線（肉眼看不到），所以看起來黑。
\item \textbf{高溫}：放出的輻射進入可見光，於是發紅 $\rightarrow$ 橙 $\rightarrow$ 黃 $\rightarrow$ 白。所以爐火、燈絲、星星、太陽都可近似看成「正在發光的黑體」——太陽其實就是一個約 5800 K 的黑體。
\end{itemize}

\textbf{注意：}「黑體一定是黑色的」是錯的。黑體指的是「\textbf{完全吸收}」的理想物體；它一旦夠熱就會發出耀眼的光。一顆星星是很好的黑體，卻一點也不黑。

現實中最好的近似，是\textbf{一個中空的盒子（空腔），上面開一個小孔}：光從小孔進去後在裡面不斷反射，幾乎逃不出來，所以那個小孔看起來全黑；而從小孔「漏」出來的輻射，就是黑體輻射。物理學家偏要研究它，是因為黑體輻射有個驚人的普適性——\textbf{它的光譜只跟溫度有關，完全不管材質、形狀、顏色}。這種「乾淨、與細節無關」的現象，正是檢驗理論的最佳試金石。
\end{補充框}

結果古典理論（瑞立–京斯定律）預測：\textbf{波長越短，輻射越強，到紫外區趨於無限大}。這顯然荒謬——一個火爐怎麼可能放出無限的能量？這個破綻史稱\textbf{紫外災難（ultraviolet catastrophe）}。

要看清古典理論「為什麼必然算出無限大」，得先認識兩個關鍵名詞：\textbf{振動模式}與\textbf{能量均分定理}。

\begin{補充框}{關鍵名詞一：振動模式（mode）}
把黑體想成一個\textbf{空腔}，裡面充滿電磁波。就像吉他弦兩端固定、只能振出特定的駐波（基音、泛音\dots），空腔裡的電磁波也只能是\textbf{剛好塞得進盒子的駐波}。每一種這樣的駐波圖樣，就叫一個\textbf{振動模式}。以一維、長度 $L$ 的盒子為例，允許的波長是
$$\lambda=\frac{2L}{n}\qquad(n=1,2,3,\dots),$$
也就是「盒子裡剛好裝整數個半波長」。

\textbf{注意：}重點不是「波長越短越塞得進去」，而是\textbf{在固定的一小段波長範圍裡，有幾個允許的模式}。這些允許值在長波長端隔得很開、在短波長端\textbf{擠得超密}；再加上真實空腔是\textbf{三維}的，同一個波長還能朝 $x$、$y$、$z$ 各種方向振盪。兩者相乘的結果是：\textbf{波長越短，單位波長內的模式數越多，而且沒有上限}——這才是「模式爆炸性地多」的真正意思。
\end{補充框}

\fig{q-modes}{0.92}{圖 6-1：左——空腔裡只有「整數個半波長」的駐波（振動模式）能存在；右——波長越短（頻率越高），模式總數急速暴增且無上限。}

\begin{補充框}{關鍵名詞二：能量均分定理（equipartition theorem）}
這是古典統計力學的一條定理：\textbf{在溫度 $T$ 的熱平衡下，能量會「公平地」分給每一個可用的模式，平均每個模式分到約 $kT$ 的能量}（$k$ 為波茲曼常數）。白話說，古典物理認為「大家有福同享」，不管哪個模式，平均都領到一樣多的能量 $kT$。
\end{補充框}

\begin{補充框}{關鍵名詞三：$kT$ 是什麼？溫度背後的「熱能尺度」}
$kT$ 是「把溫度換算成能量」的尺度——某個溫度下，熱擾動「典型」給每個模式的能量。
\begin{itemize}\setlength{\itemsep}{1pt}
\item $k$＝\textbf{波茲曼常數（Boltzmann constant）}$\approx 1.38\times10^{-23}$ J/K，是「溫度 $\leftrightarrow$ 能量」的換算係數。
\item $T$＝\textbf{絕對溫度}（克氏溫標 K，從絕對零度 $-273\,^\circ$C 起算）。一定要用絕對溫度，不能代攝氏度數。
\end{itemize}
室溫（$T\approx300$ K）時
$$kT \approx 1.38\times10^{-23}\times300 \approx 4\times10^{-21}\ \text{J} \approx 0.025\ \text{eV}\ \Big(\approx\tfrac{1}{40}\text{ eV}\Big).$$
記住這把尺「\textbf{室溫 $kT\approx0.025$ eV}」，之後判斷量子效應顯不顯著很方便。注意 $kT$ \textbf{不是}某顆粒子的確切能量，而是\textbf{典型尺度／平均值}。
\end{補充框}

把兩件事乘起來，災難就出現了：
$$\text{總輻射能} \;\approx\; (\text{模式數目}) \times (\text{每個模式 } kT).$$
短波長的模式\textbf{數目沒有上限}，每個又\textbf{照領 $kT$}，於是短波長（紫外）累積的能量 $\rightarrow\infty$。這不是哪個公式算錯，而是「能量連續、人人有份 $kT$」這套古典假設，邏輯上\textbf{必然}導致無限大。

\fig{q-blackbody}{0.80}{圖 6-2：黑體輻射光譜。實線是符合實驗的量子論曲線（有峰值、會收斂）；灰色虛線是古典理論——波長越短衝得越高、發散成「紫外災難」。}

\subsection*{\sffamily B.\ 普朗克的量子假設（1900）}

普朗克（Planck）發現：只要假設物體吸收或放出能量時\textbf{不是連續的，而是一份一份}，每份大小為
$$\boxed{\,E = h\nu\,}$$
（$\nu$ 為頻率，$h=6.63\times10^{-34}$ J·s 為\textbf{普朗克常數（Planck constant）}），算出的曲線就\textbf{完美符合實驗}。為什麼這樣就能化解災難？因為\textbf{高頻（短波長）模式的「一份」$h\nu$ 非常大}，在溫度 $T$ 下根本湊不出一份能量（$h\nu\gg kT$），於是這些模式被「凍結」、領不到能量、不參與輻射——無限大就消失了。這個「能量量子化」的假設，是整個量子論的起點。

由於 $h$ 極小（$10^{-34}$ 等級），能量的「顆粒」在日常尺度上細到看不出來，所以我們才會以為能量是連續的。\textbf{量子效應只在微觀才顯著}——這條伏筆會貫穿全章。

\begin{延伸框}{普朗克的「絕望之舉」——連發明者都不相信}
課本把量子化寫成漂亮的解題，容易給人「天才靈光一閃」的錯覺。真相更有人味：普朗克\textbf{並不相信}能量真的量子化。他稱這個假設是「一個絕望之舉（an act of desperation）」，把 $h$ 當成純粹的數學湊合，在之後十多年裡不斷嘗試把 $h\to0$、回到連續的古典圖像，想證明量子化只是表象——他\textbf{失敗了}。

真正讓「能量量子」被當成物理實在的，是它在\textbf{別的、看似無關的現象}上也成功：愛因斯坦（1905）用同一個 $h\nu$ 解釋光電效應（見 6-2），之後比熱、原子光譜也一個個倒向量子。這是科學史的好教材：一個正確的觀念，常常連提出者都半信半疑，要靠它在多個獨立現象上都奏效，才被接受。

至於更深的一層——「\textbf{為什麼我們的宇宙是量子的}？$h$ 為何存在、能不能由更基本的原理導出？」——至今仍是物理學的\textbf{開放問題，目前無定論}。量子力學「描述」得極其成功，卻沒告訴我們它「為什麼」長這樣。
\end{延伸框}

\section*{\sffamily 6-2\quad 光的粒子性：光電效應}

\subsection*{\sffamily A.\ 實驗事實}

紫外光照在金屬表面，會打出電子（光電子）。古典波動說認為「光夠強、照夠久就能打出電子」，但實驗結果處處與它衝突：

\begin{center}
\renewcommand{\arraystretch}{1.35}
\begin{tabular}{p{0.46\linewidth} p{0.40\linewidth}}
\hline
\textbf{觀察到的事實} & \textbf{古典波動說的預測（錯）} \\
\hline
低於某\textbf{遏止頻率 $\nu_0$}，再亮也打不出電子 & 只要夠亮就一定打得出 \\
高於 $\nu_0$，\textbf{瞬間}就有電子，毫無延遲 & 弱光要「累積」一段時間才行 \\
電子最大動能只跟\textbf{頻率}有關，與亮度無關 & 動能應該隨光強（亮度）增加 \\
光越亮 $\rightarrow$ 電子\textbf{數目}越多（但動能不變） & （這一點與光強有關，倒是符合） \\
\hline
\end{tabular}
\end{center}

\subsection*{\sffamily B.\ 愛因斯坦的光量子（1905）}

愛因斯坦主張：光本身就是一顆顆\textbf{光子（photon）}，每顆能量 $E=h\nu$。一顆光子把能量\textbf{整份}交給一個電子；電子要先付出\textbf{功函數 $W$}（逃出金屬的最低能量），剩下的才變成動能：
$$\boxed{\,K_{\max}=h\nu-W\,},\qquad W=h\nu_0.$$

這一口氣解釋了上表所有矛盾：頻率太低（$h\nu<W$），一顆光子能量不夠，再多顆也沒用；能量是「一對一」整份傳遞，所以瞬時、且動能只看頻率不看亮度。

\begin{補充框}{什麼是功函數（work function）？金屬為什麼「扣住」電子？}
金屬導電，是因為內部有大量可以到處跑的傳導電子。但「能在金屬\textbf{裡面}自由移動」不等於「能自由\textbf{離開}金屬」——就像魚能在魚缸裡自由游，卻跳不出魚缸。

金屬表面有一道「能量牆」：電子一旦離開表面，背後會留下淨正電荷把它往回拉。整體效果是，金屬內部像一個\textbf{位能井（凹槽）}，電子待在井底，要離開就得爬出這道牆。

\textbf{功函數 $W$＝把最容易跑的電子剛好帶到金屬外（速度為零）所需的最小能量＝這道牆的高度。}可以想成一個\textbf{碗}：電子是碗裡的水，能在碗裡自由晃動（＝導電），但要潑出碗外，必須先給它足夠能量爬過碗緣——碗緣的高度就是功函數。

這也解釋了\textbf{遏止頻率 $\nu_0$}：一顆光子能量 $h\nu$ 必須先「付清」這道牆 $W$，才剩得下動能，所以臨界條件是 $h\nu_0=W$。\textbf{注意：}光「夠亮」不代表打得出電子——亮只是光子\textbf{數目}多；每顆光子能量 $h\nu$ 若不足 $W$，每一顆都翻不過牆。關鍵在\textbf{單顆光子的能量（頻率）}，不是總亮度。
\end{補充框}

\fig{q-photoelectric}{0.74}{圖 6-3：光電子最大動能 $K_{\max}$ 對入射光頻率 $\nu$ 的關係。直線斜率就是普朗克常數 $h$，與 $y$ 軸截距對應 $-W$；橫軸截距即遏止頻率 $\nu_0$。功函數大的金屬（B）需要更高頻的光才打得出電子。}

光電效應就在我們身邊：太陽能電池、自動門與樓梯的光感測、相機測光、影印機\dots\ 都用到光電／光生伏特效應。

\section*{\sffamily 6-3\quad 物質的波動性：德布羅意波}

光既然是波又是粒子，德布羅意（de Broglie，1924）大膽反過來問：那\textbf{物質（電子）會不會也是波}？他提出，任何動量為 $p$ 的粒子都伴隨一個波，波長為
$$\boxed{\,\lambda=\frac{h}{p}=\frac{h}{mv}\,}$$
這叫\textbf{物質波}，或\textbf{德布羅意波（de Broglie wave）}。

\begin{補充框}{這條 $\lambda=h/p$ 是怎麼來的？不是憑空猜的}
它是「推理＋對稱性」的產物，靠三步：

\textbf{第一步——先看「光」。}光是波（有波長 $\lambda$、頻率 $\nu$），但量子論說它也是一顆顆光子，能量 $E=h\nu=hc/\lambda$。再由相對論，光子雖無質量卻有\textbf{動量（momentum）}，且 $p=E/c$。把兩者合起來：
$$p=\frac{E}{c}=\frac{hc/\lambda}{c}=\frac{h}{\lambda}\;\Longrightarrow\;\lambda=\frac{h}{p}.$$
也就是說，對「光」而言，波長和動量\textbf{早就}被 $\lambda=h/p$ 綁在一起了。

\textbf{第二步——大膽的對稱。}德布羅意想：既然一向被當成「純粹的波」的光竟帶有粒子的動量；那一向被當成「純粹的粒子」的電子，會不會也伴隨一個波？\textbf{自然界應該是對稱的}，於是他把同一條關係反過來套到物質上。這不是亂猜，是「光能這樣，物質沒理由不能」的論證。

\textbf{第三步——它還順手解釋了玻爾的謎團。}若電子是繞核的\textbf{波}，只有「軌道剛好裝得下整數個波長」的駐波才穩定：$2\pi r=n\lambda$，整理後正好得到玻爾硬性規定的量子化條件 $mvr=n\hbar$。物質波因此一口氣解釋了「為什麼能階是分立的」。
\end{補充框}

不久，戴維森–革末（1927）用電子打鎳晶體，看到了和波一模一樣的\textbf{繞射圖樣}，證實電子真的有波長。\textbf{電子顯微鏡}正是利用電子波長極短，能看清比光學顯微鏡更小的東西。

那為什麼棒球、人不會「繞射」？因為 $h$ 極小，宏觀物體動量 $p$ 巨大，算出來的 $\lambda$ 小到失去意義：一顆棒球的德布羅意波長約 $10^{-34}$ m，比原子核還小 $10^{19}$ 倍。波動性（繞射、干涉）只有在波遇到「與波長尺度相當」的狹縫或障礙時才明顯——電子波長（約 $0.1$ nm）和原子間距相當，用晶體當光柵就看得到；但\textbf{宇宙中根本沒有}小到能讓棒球繞射的狹縫。

\begin{延伸框}{萬物皆波，為何我們感覺不到自己的波？退相干與前沿實驗}
說「宏觀物體沒有波動性」其實不精確。更正確的說法是：\textbf{有，但波長與「相干時間」都小到永遠測不到}。除了波長太小，還有一個更關鍵的原因——\textbf{退相干（decoherence）}。

要表現出波動性，各條路徑之間必須維持穩定的相位關係（相干性）。但宏觀物體無時無刻都在和環境（空氣分子、光子、熱輻射）發生無數次交互作用，每一次都把它的量子相位資訊「洩漏」給環境，使不同路徑間的干涉\textbf{極快地、不可逆地}被抹平。一顆塵埃只要被幾個空氣分子或光子撞一下，相干性就在遠短於 $10^{-30}$ 秒內消失。所以宏觀的量子疊加不是「不存在」，而是\textbf{存在的時間短到絕無可能觀測}。這是目前「量子 $\rightarrow$ 古典」過渡的主流解釋。

而這條邊界至今仍在被推進：實驗物理學家不斷讓\textbf{越來越大的物體}展現量子干涉——從電子、原子，到 C$_{60}$ 富勒烯，再到上萬個原子質量的大分子，都做出了雙狹縫干涉。「量子疊加的尺寸\textbf{有沒有一個基本的上限}？」仍是\textbf{開放問題，目前無定論}。
\end{延伸框}

\section*{\sffamily 6-4\quad 波粒二象性與雙狹縫}

把電子\textbf{一顆一顆}發射穿過雙狹縫，打在屏幕上。每顆電子只留下\textbf{一個點}（粒子性）；但累積幾萬顆後，這些點竟組成\textbf{干涉條紋}（波動性）——彷彿每顆電子都「同時通過兩條縫、自己和自己干涉」。

\fig{q-doubleslit}{0.86}{圖 6-4：單電子雙狹縫實驗。電子一個一個打（每個只是一個點），數目越多，越清楚地累積出干涉條紋。}

\begin{補充框}{「同時通過兩條縫」是什麼意思？這就是量子「疊加」}
\textbf{疊加（superposition）}＝在被測量之前，量子系統可以同時處於多個狀態的「組合」。電子的狀態可以是「走左縫」和「走右縫」\textbf{同時}的疊加，測量時才隨機得到其中一個確定結果。

\textbf{注意——這是最該破除的誤解：}疊加\textbf{不是}「它其實走了某一邊、只是我們不知道」。後者是古典機率的想法，是錯的。兩者可以用實驗分辨：
\begin{itemize}\setlength{\itemsep}{1pt}
\item 若電子「其實走了某一邊、只是我們不知道」（古典），把很多次結果疊起來，應該是兩條縫各自的亮區相加，\textbf{沒有條紋}。
\item 但實驗看到\textbf{干涉條紋}——只有當電子「兩條路同時存在、彼此干涉」（真正的疊加）才會發生。
\end{itemize}
所以\textbf{干涉條紋就是「疊加是真的、不是無知」的鐵證}。（薛丁格那隻「又死又活的貓」正是把這個疊加放大到宏觀，用來凸顯荒謬、質疑理論，並不是說貓真的又死又活。）
\end{補充框}

最關鍵也最詭異的一步是：若你在縫口裝偵測器、想看電子到底走哪條縫——\textbf{干涉條紋會消失}，電子立刻表現得像純粹的粒子。「看」這個動作本身改變了結果。這就是\textbf{波耳互補原理（complementarity principle）}：波動性與粒子性是互補的兩面，無法同時觀察到。

\begin{延伸框}{「觀測」為什麼會破壞干涉？量子測量問題（至今無共識）}
互補原理只\textbf{描述}現象，沒有\textbf{解釋}「觀測為什麼會這樣」。背後藏著物理學至今爭論不休的\textbf{量子測量問題}。

量子力學說：測量前電子處於各種可能性的疊加，由\textbf{波函數（wavefunction）}$\psi$ 描述，平滑地、決定論地演化；但我們每次測量只會得到一個明確結果。從「疊加」到「單一結果」這一步——所謂\textbf{波函數塌縮（wavefunction collapse）}——是瞬間、隨機、不可逆的，並不在原本的方程式裡，而是額外硬加上去的規則。同一個理論裡塞了兩套互不相容的演化規則，而「什麼時候算一次測量」卻沒有客觀定義。

主流物理學家對「世界實際發生了什麼」各有立場、互不相讓（哥本哈根詮釋、多世界詮釋、退相干、客觀塌縮理論、隱變數／玻姆力學\dots）。它們對「實驗預測」幾乎完全一致，對「真相」卻南轅北轍。即使 2022 年諾貝爾獎肯定的貝爾不等式實驗已排除「定域隱變數」，大多數立場仍然存活。\textbf{到目前為止，沒有任何一種詮釋被公認為正解——這是開放問題，目前無定論。}

\textbf{要破除的迷思：}常聽到「是\textbf{意識}讓波函數塌縮」。主流並\textbf{不}認為需要意識——一個偵測器、甚至只是與環境的交互作用（退相干）就夠了。「意識造成塌縮」是少數且有爭議的立場，不是定論。
\end{延伸框}

\section*{\sffamily 6-5\quad 原子光譜與能階}

\subsection*{\sffamily A.\ 線狀光譜之謎}

把氫氣放電發出的光用稜鏡分開，\textbf{不是連續彩虹，而是幾條固定波長的亮線}（\textbf{線狀光譜，line spectrum}）。每種元素的譜線像「指紋」一樣獨一無二。古典物理完全無法解釋：為什麼只有這幾個特定波長？

\subsection*{\sffamily B.\ 玻爾模型與能階（1913）}

玻爾（Bohr）假設：電子只能待在\textbf{特定的能階（energy level）}上，每個能階有固定能量
$$E_n=-\frac{13.6}{n^2}\ \text{eV}\qquad(n=1,2,3,\dots).$$
在能階上電子\textbf{不輻射}；只有當電子從高能階 $E_i$ \textbf{躍遷（transition）}到低能階 $E_f$ 時，才放出一顆光子，能量恰為兩能階之差：
$$\boxed{\,h\nu = E_i - E_f\,}.$$
因為能階是分立的，能量差也是分立的，所以放出的光只有幾個特定頻率——\textbf{線狀光譜就此解釋}。

\begin{補充框}{為什麼能量前面是「負號」？「束縛態」的意思}
能量的「絕對數值」沒有意義，\textbf{只有能量差有意義}，所以我們可以自由選擇哪裡當零點。原子物理裡大家約定：\textbf{零能量＝電子剛好脫離原子核、跑到無限遠、且速度為零}（剛剛好被游離的瞬間）。

選定這個零點後，負號就說得通了：電子和原子核\textbf{互相吸引}，要把被吸住的電子拉到無限遠，你必須\textbf{對它做功、輸入能量}。既然把它送到「零」還需要額外加能量，代表它原本的能量\textbf{比零還低}，也就是\textbf{負的}。這種「被綁住、能量低於自由狀態」的狀態，就叫\textbf{束縛態（bound state）}。原子裡的電子全是束縛態，所以能階全為負：
$$E_n<0\ (\text{被綁住}) \quad<\quad E=0\ (\text{剛好自由}) \quad<\quad E>0\ (\text{自由且有動能}).$$

兩個馬上有用的結論：
\begin{itemize}\setlength{\itemsep}{1pt}
\item \textbf{越負＝綁得越緊}。$E_1=-13.6$ eV 最負，所以基態電子被綁得最牢；$n$ 越大越接近 0，綁得越鬆，到 $n=\infty$ 時 $E=0$（快要自由）。
\item \textbf{能階的絕對值＝游離能（ionization energy）}。$|E_1|=13.6$ eV 就是「把基態電子完全扯離原子，至少要給它 13.6 eV」。
\end{itemize}
這跟「把零點定在無限遠的重力位能」一模一樣：一顆被地球綁住的衛星總能量是負的，要讓它逃逸就得加能量把它推到「能量＝0」。\textbf{負能量＝被束縛、跑不掉}，不是什麼神祕的東西。
\end{補充框}

\fig{q-levels}{0.66}{圖 6-5：氫原子能階圖（$E_n=-13.6/n^2$ eV）。電子由高能階躍遷到低能階時放出光子；落到 $n=1$ 的是來曼系（紫外），落到 $n=2$ 的是巴耳末系（可見光）。能階越往上越密，到 $n=\infty$（$E=0$）即游離。}

玻爾的「電子只能待在特定軌道、且不輻射」\textbf{完全違反}古典電磁學——古典上，繞圈的電子會持續輻射、瞬間墜入原子核。這正說明：\textbf{牛頓力學與古典電磁學在原子尺度下並不適用}，必須由量子力學取代。

\begin{延伸框}{為什麼電子不會一邊繞一邊掉進原子核？}
古典電磁學有個鐵則：\textbf{加速中的電荷會輻射電磁波}（這正是天線發射無線電的原理）。電子繞核是向心加速運動，所以古典上它\textbf{必須}持續輻射能量。把數字算出來會嚇一跳：一個古典氫原子的電子會在約 $10^{-11}$ 秒內把能量輻射殆盡、螺旋墜入原子核。換句話說，\textbf{如果古典物理是對的，所有原子應該在一兆分之一秒內崩塌，物質根本不可能存在}。我們的存在本身，就是古典物理失敗的證據。

真正的解釋來自量子力學，核心是\textbf{不確定性原理（uncertainty principle）}：你無法同時把一個粒子的位置和動量都測得任意準，
$$\Delta x\,\Delta p \gtrsim \frac{\hbar}{2}.$$
\textbf{這不是儀器不夠精密}，而是粒子的本性（波動性的必然結果）：想讓位置很確定（波包很窄），就得混入很多種波長，動量反而很不準；反之亦然。若電子想「掉進原子核」，等於把位置 $\Delta x$ 壓到極小，動量的不確定 $\Delta p$ 就被迫變得極大，動能 $\sim(\Delta p)^2/2m$ 跟著暴增、把它再撐開。「被核吸引想往內掉」與「位置壓縮導致動能暴增想往外撐」兩者平衡，得到一個\textbf{能量最低的穩定基態}，對應有限大小的原子。

\textbf{一句話：}原子不塌縮，不是因為電子「繞得夠快」，而是因為量子力學規定能量有個下限，再低就違反不確定性原理——基態之下沒有更低的階可掉了。「你能站在地上、不穿過地板」，靠的就是這件事。（單一原子的穩定性\textbf{已由量子力學完整解釋，是定論}，這一點與 6-4 那個仍無定論的測量問題不同。）

\textbf{注意：}「電子像行星繞太陽」只是方便的舊圖像，並不正確。量子力學裡電子是分布在空間的\textbf{機率雲（軌域）}，沒有確定的軌道，問「它確切在哪」本身就不是好問題。
\end{延伸框}

\section*{\sffamily 6-6\quad 本章重點整理}

\begin{補充框}{一頁複習（公式與關鍵結論）}
\begin{itemize}\setlength{\itemsep}{2pt}
\item \textbf{量子化}：能量一份一份，$E=h\nu$，$h=6.63\times10^{-34}$ J·s。起於黑體輻射（化解紫外災難）。
\item \textbf{光電效應}：光是光子，$K_{\max}=h\nu-W$，$W=h\nu_0$。證明光的\textbf{粒子性}。功函數 $W$＝電子逃出金屬的「過路費」。
\item \textbf{德布羅意波}：$\lambda=h/p=h/mv$。物質也有\textbf{波動性}；因 $h$ 太小，宏觀物體波長極小、又因退相干，故測不到。
\item \textbf{波粒二象性／雙狹縫}：電子一顆顆打仍累積出干涉（疊加是真的）；一旦「觀測路徑」干涉就消失（互補原理）。
\item \textbf{能階與光譜}：$E_n=-13.6/n^2$ eV（負號＝束縛態）；躍遷放光子 $h\nu=E_i-E_f$，解釋線狀光譜。原子尺度下古典力學失效，由不確定性原理撐住原子不塌縮。
\end{itemize}
\textbf{常用近似}：$E(\text{eV})=1240/\lambda(\text{nm})$；室溫 $kT\approx0.025$ eV $\approx\tfrac{1}{40}$ eV。
\end{補充框}

\end{document}
